\done\todo{Develop an abstract: The abstract summarizes the information contained in the rest of this document. An abstract includes:(a)the research question or engineering problem, (b) procedures used, (c) data, (d) interpretation, and (e) conclusions. It also may include any possible research applications. It should be limited to these essential elements. Since this is a proposal, the conclusion part should be a potential conclusion.}

\done\todo{Please check the abstract and ensure it is appropriate based on the whole document.}
Software security is an ongoing concern as software continues to develop; therefore, teaching and learning methods of software security becomes increasingly important. An obstacle to this is the lack of high-quality, real-world examples of vulnerabilities and their solutions to reference. In order to meet these needs, we want to propose a methodology that first finds relevant, modern, and patched vulnerabilities from the National Vulnerability Database (NVD). Then, these examples will be examined to determine their complexity and identify the prerequisite knowledge needed in order to implement said patches so that we can further make conclusions for educational implementations. This will be further supported with automated filtering methods, including the potential use of machine learning for CWE classification and scoring difficulty. The resulting dataset is expected to support the development of teaching and learning in software security curriculum by offering a curated set of documented, robust and contextualized examples.